\section{Related Work}
\label{sec:related}

\subsection{RF Fingerprinting and LoRa RFFI}
Radio frequency fingerprint identification exploits transmitter-specific hardware impairments for device verification~\cite{brik2008wireless_device_id,jian2020deeplearning_rf_fingerprinting}.
Surveys summarize datasets, signal representations, and deployment challenges for IoT-oriented RFFI~\cite{ahmed2024lora_rffi_survey,xie2021physical_layer_auth_survey}.
LoRa RFFI is complicated by chirp spread spectrum modulation and by sensitivity to collection conditions~\cite{shen2021lora_spectrogram_cnn,elmaghbub2021lora_wild}.
The OSU LoRa corpus systematically exposes cross-day, cross-location, cross-configuration, and cross-receiver degradation~\cite{hamdaoui2022deeplearning,elmaghbub2021lora_wild}.
Our work adopts the cross-receiver subset of that corpus and evaluates under strict block-disjoint calibration/query splits rather than random window mixing.

\subsection{OOB Evidence in LoRa RFFI}
WideScan and follow-on studies show that out-of-band distortion can strengthen device signatures caused by power-amplifier nonlinearity and filter roll-off~\cite{elmaghbub2020widescan,hamdaoui2022deeplearning}.
IoTJ and TIFS studies further emphasize robust fingerprint extraction under variable conditions~\cite{zhou2021robust_rffi_iotj}.
We do not claim the first use of OOB evidence.
Instead, we study what happens to OOB-guided representations when the receiver changes: OOB can become receiver-entangled rather than purely device-discriminative, which motivates diagnosis before calibration.

\subsection{Cross-Receiver and Receiver-Agnostic RFFI}
Receiver mismatch is a known open challenge in LoRa deployment studies~\cite{elmaghbub2021lora_wild}.
Receiver-agnostic training with adversarial or collaborative objectives has been proposed to reduce receiver sensitivity, often requiring multi-receiver training or coordinated data collection~\cite{shen2023receiveragnostic,ma2025receiver_independent_rffi,shen2022scalable_lora_rffi}.
Recent cross-receiver generalization work explores feature disentanglement and adversarial training to learn receiver-invariant representations, but typically under different receiver populations, signal regimes, and training protocols than our single indoor RX1/RX2 setup.
Channel-robust and contrastive LoRa RFFI methods improve embedding compactness under channel variation but typically retrain or adapt the backbone under different enrollment assumptions~\cite{shen2022scalable_lora_rffi,ma2025receiver_independent_rffi}.
\textbf{Direct numerical comparison with these systems is not apples-to-apples} because they differ in receiver sets, devices, source-data access, target-label availability, and calibration protocols.
Our experiments instead compare against representative \emph{same-protocol} baselines under identical block-disjoint splits (Table~\ref{tab:sota_baselines}).

\subsection{Source-Free and Unsupervised Adaptation}
Test-time adaptation methods such as entropy minimization update normalization or classifier statistics using unlabeled target data~\cite{wang2021tent}.
Source-free cross-receiver RFFI adapts pre-trained models with contrastive or pseudo-label objectives on target-receiver traffic without source data or target labels.
These approaches target a different deployment mode from ours: we assume a frozen same-receiver backbone and a small set of labeled target-receiver calibration windows.
Our TTA negative baselines, threshold sweep, and feature-alignment baselines show that unsupervised repair remains near chance under severe RX1$\rightarrow$RX2 collapse, supporting labeled calibration as the more reliable option under our protocol.

\subsection{Few-Shot and Prototype Classification}
Prototypical classification represents each class by the mean embedding of a support set and assigns queries by nearest-prototype distance~\cite{snell2017prototypical}.
Few-shot RFFI and open-set studies often use episodic meta-learning or multiple enrollment captures per device.
We do not claim a new prototype algorithm.
Instead, we compare RCPA-T against K-shot logistic regression and classifier-head fine-tuning under a LoRa-specific protocol with one file per device per receiver, where $K$ denotes labeled windows and calibration/support/query blocks are explicitly disjoint.
Same-protocol results (Table~\ref{tab:sota_baselines}) show that RCPA-T is competitive with linear probe while avoiding gradient-based target adaptation.

\subsection{Positioning of This Work}
Table~\ref{tab:related_compare} summarizes how RCPA differs from representative directions.
Prior backbone work~\cite{paper1_placeholder} studies OOB-guided architecture design and source-only same-receiver cross-day robustness; the present paper explains and mitigates the cross-receiver limitation reported there through diagnosis and target-receiver calibration on a frozen model.

\begin{table}[t]
\caption{Positioning against representative cross-receiver RFFI directions.}
\label{tab:related_compare}
\centering
\small
\begin{tabular}{p{0.22\linewidth}p{0.68\linewidth}}
\toprule
Direction & vs.\ RCPA (this work) \\
\midrule
Receiver-agnostic train & Multi-RX retraining; we freeze backbone and calibrate with $K$ target windows \\
Source-free cross-RX & Unlabeled target only; we show TTA/pseudo insufficient under collapse \\
Feature disentangle & Retrain/adapt backbone; we use post-hoc diagnosis + prototypes \\
Few-shot prototype RFFI & Generic episodic few-shot; we add OOB entanglement diagnosis + block-disjoint LoRa protocol \\
Prior OOB hybrid backbone~\cite{paper1_placeholder} & Architecture + cross-day; we diagnose cross-RX failure + RCPA-T calibration \\
\bottomrule
\end{tabular}
\end{table}
